% ---
% RESUMOS
% ---

% RESUMO em português
\setlength{\absparsep}{18pt} % ajusta o espaçamento dos parágrafos do resumo
\begin{resumo}[Resumo]

 Dados de metabolômica não direcionada obtidos por cromatografia líquida acoplada à espectrometria de massas (LC-MS) chegam do processamento primário com erros sistemáticos do equipamento, e as etapas seguintes de refino — correção do sinal, filtragem, tratamento de valores ausentes, normalização, transformação e escalonamento — costumam ser feitas em ferramentas separadas, com escolhas manuais baseadas majoritariamente no conhecimento dos pesquisadores. Este trabalho apresenta o MetaRefino, aplicativo para Windows que assiste o pesquisador nesse refino a partir das matrizes exportadas pelo Progenesis QI: valida a entrada, diagnostica a forma da deriva instrumental, propõe as fronteiras de lote e o método de correção com justificativa, aplica as correções LOESS e SERRF ou mantém o sinal sem correção, executa o pré-processamento e as análises estatísticas univariadas e multivariadas, e exporta a matriz no formato do MetaboAnalyst. A escolha do corretor por um modelo de aprendizado de máquina não superou uma regra fixa e deu lugar, na interface, a uma medição sobre a própria matriz. A ferramenta foi avaliada com os dados de um único estudo do laboratório parceiro, com matrizes sintéticas derivadas deles e com a devolutiva escrita de seus usuários. O principal resultado veio da comparação entre réguas de qualidade: medida nos próprios controles de qualidade usados no ajuste, a correção LOESS fica à frente da ausência de correção nos quatro modos cromatográficos; medida fora da amostra, fica atrás dela em três. Conclui-se que a ferramenta integra e torna auditáveis as etapas do refino quimiométrico, com a decisão final mantida com o pesquisador; a generalização para outros estudos, laboratórios e equipamentos permanece por avaliar.
\vspace{\onelineskip}

 \noindent
 \textbf{Palavras-chave}: metabolômica; LC-MS; correção de deriva instrumental; controle de qualidade; quimiometria.
\end{resumo}

% ABSTRACT in english
\begin{resumo}[Abstract]
 \begin{otherlanguage*}{english}
   Untargeted metabolomics data acquired by liquid chromatography coupled to mass spectrometry (LC-MS) leave primary processing with systematic instrument errors, and the subsequent refinement steps --- signal correction, filtering, missing value treatment, normalization, transformation and scaling --- are usually carried out in separate tools, with manual choices based mainly on the researchers' knowledge. This work presents MetaRefino, a Windows application that assists researchers in this refinement, starting from the matrices exported by Progenesis QI: it validates the input, diagnoses the shape of the instrumental drift, proposes batch boundaries and a correction method with a justification, applies LOESS and SERRF correction or keeps the signal uncorrected, runs the preprocessing and the univariate and multivariate statistical analyses, and exports the matrix in the MetaboAnalyst format. Corrector selection by a machine learning model did not outperform a fixed rule and was replaced, in the interface, by a measurement on the matrix itself. The tool was evaluated on data from a single study of the partner laboratory, on synthetic matrices derived from them, and through written feedback from its users. The main result came from comparing quality yardsticks: measured on the same quality control samples used for fitting, LOESS correction outperforms no correction in all four chromatographic modes; measured out of sample, it falls behind in three. The tool integrates the chemometric refinement steps and makes them auditable, while the final decision remains with the researcher; generalization to other studies, laboratories and instruments remains to be evaluated.

   \vspace{\onelineskip}

   \noindent
   \textbf{Keywords}: metabolomics; LC-MS; instrumental drift correction; quality control; chemometrics.
 \end{otherlanguage*}
\end{resumo}